% Generated by scripts/assemble.py; corrections belong in source/corrections.json.
\examhead{175}{西南财经大学 2026 年数学分析试卷}{165}
% p166-b11; source PDF page 166
\blockanchor{p166-b11}
\question{1}求极限  \displayformula{\lim_{n\to\infty}\left(1-\frac{1}{2^{2}+2}\right)^{\frac{1}{n}}\left(1-\frac{1}{3^{2}+3}\right)^{\frac{1}{n}}\cdots\left(1-\frac{1}{n^{2}+n}\right)^{\frac{1}{n}}}\par
% p166-b12; source PDF page 166
\par\Needspace{6\baselineskip}
\blockanchor{p166-b12}
\question{2}求 $z = \frac{1}{2} (x^4 + y^4)$ 在条件 $x + y = a(x, y \geq 0)$ 下的最小值, 并在此基础上证明不等式:\par
\nopagebreak[4]
% p166-b13; source PDF page 166
\blockanchor{p166-b13}
\displayformula{\frac {x ^ {4} + y ^ {4}}{2} \geq \left(\frac {x + y}{2}\right) ^ {4}.}
% p166-b14; source PDF page 166
\blockanchor{p166-b14}
\question{3}已知 $4\frac{\partial^2f}{\partial u^2} - \frac{\partial^2f}{\partial v^2} = 1$，二元函数 $g(x, y) = f(x^2 + y^2, xy)$，证明：$\frac{\partial^2g}{\partial x^2} - \frac{\partial^2g}{\partial y^2} = x^2 - y^2$.\par
% p166-b15; source PDF page 166
\blockanchor{p166-b15}
\question{4}讨论级数 $\sum_{n=1}^{\infty} \frac{(-1)^n}{n\sqrt[n]{n}}$ 的敛散性, 如果收敛, 是绝对收敛还是条件收敛?\par
% p166-b16; source PDF page 166
\par\Needspace{6\baselineskip}
\blockanchor{p166-b16}
\question{5}已知二元函数\par
\nopagebreak[4]
% p166-b17; source PDF page 166
\blockanchor{p166-b17}
\displayformula{f (x, y) = \left\{ \begin{array}{l l} e ^ {- (x + y)}, & x y > 0; \\ 1, & x y \leq 0. \end{array} \right.}
% p166-b18; source PDF page 166
\blockanchor{p166-b18}
积分区域 $D = \{(x,y)\mid 0\leq |x| + |y|\leq 1\}$ ，求 $\iint_{D}f(x,y)\mathrm{d}x\mathrm{d}y.$\par
% p166-b19; source PDF page 166
\blockanchor{p166-b19}
\question{6}已知  $u_{1}>0$ ,  $u_{n}e^{u_{n+1}}=e^{u_{n}}-1(n=1,2,\cdots)$ , 证明:  $\lim_{n\to\infty}u_{n}$  存在, 并求其值.\par
% p166-b20; source PDF page 166
\par\Needspace{6\baselineskip}
\blockanchor{p166-b20}
\question{7}已知 $f(x)$ 二阶连续可微, 证明:\par
\nopagebreak[4]
% p166-b21; source PDF page 166
\blockanchor{p166-b21}
\displayformula{\int_ {a} ^ {b} f (x) \mathrm{d} x - \frac {1}{2} (b - a) [ f (a) + f (b) ] = \frac {1}{2} \int_ {a} ^ {b} f ^ {\prime \prime} (x) (x - a) (x - b) \mathrm{d} x.}
% p166-b22; source PDF page 166
\blockanchor{p166-b22}
\question{8}设函数 $f$ 在 $[0, +\infty)$ 上连续可导, $f(0) > 0$ 且 $f'(x) \geq 0$, $x \in [0, +\infty)$, 证明: 若 $\int_0^{+\infty} \frac{1}{f(x) + f'(x)} \mathrm{d}x$ 收敛, 则反常积分 $\int_0^{+\infty} \frac{1}{f(x)} \mathrm{d}x$ 也收敛.\par
